\documentclass[10pt,a4paper]{amsart}
\usepackage[margin=19mm]{geometry}
\usepackage{amsmath,amssymb,mathtools}
\usepackage[scr=rsfs,cal=euler]{mathalfa}
\usepackage{xfrac}
\usepackage[unicode,hidelinks,bookmarks=false]{hyperref}
\newcommand{\ie}{i.e.\ }
\newcommand{\cf}{cf.\ }
\newcommand{\resp}{respectively\ }
\newcommand{\Subsection}{\subsection}
\providecommand{\nobreakdash}{\nobreak\hbox{-}}
\setlength{\emergencystretch}{2em}
\multlinegap0pt
\usepackage{amssymb}
\usepackage{mathtools}
\newtheorem{theorem}{Theorem}
\title{Tighter Bounds for Algorithmic Complexity Estimation Using a Reusable Code-Based Block Decomposition Method}
\author{Eduardo Yuji Sakabe, Felipe S. Abrah\~ao, Santiago Hern\'andez-Orozco, Ricardo Gudwin, Hector Zenil}
\date{Source: arXiv:2606.23471v1 (CC BY 4.0)}
\begin{document}
\maketitle

\noindent\emph{Source and licence.} Tighter Bounds for Algorithmic Complexity Estimation Using a Reusable Code-Based Block Decomposition Method. Version arXiv:2606.23471v1 (CC BY 4.0), distributed by arXiv under the Creative Commons Attribution 4.0 International licence, http://creativecommons.org/licenses/by/4.0/, as recorded in the arXiv licence metadata for this exact version and shown on the abstract page https://arxiv.org/abs/2606.23471v1. Full licence notice: \texttt{licenses/arxiv-2606.23471-CC-BY-4.0.txt}.\par\smallskip

\noindent\textit{Editorial scope.} Counted unit: the article's theorem thm:bdm2-reuse-improvement (source0.tex lines 470--554). The article's complete original proof of it is delivered with it (source0.tex lines 556--683, one \texttt{proof} environment of 128 lines). No mathematics is written, repaired or re-derived: the statement and the proof are the article's own lines, in the article's own notation, and no retained line is substituted or re-flowed.  Retained uncounted as the source-local context the proof needs: lines 376--395; lines 407--414.  Nothing was omitted from any retained span, so the package carries no omission record.  No retained line contains a citation, so the wrapper carries no bibliography and no bibliography entry.  No retained line contains a figure environment, an image-inclusion command or drawing code, so no image is delivered and the package declares no asset dependency.  Editorial formatting only, in addition: an AMS \texttt{amsart} preamble that restates the article's own declarations and macros used by the retained lines, namely the environments \texttt{theorem} on separate counters, together with the standard packages those lines need.  The retained environments print in this document as Theorem 1 in order of appearance, whereas the article prints the same items with the numbering of its own full document.  The complete native source of the article as distributed in the arXiv e-print for this exact version is bundled unchanged as \texttt{sources/203435/source0.tex}.
\begin{theorem}[Independent-explanation upper bound]
\label{thm:bdm2-independent-limit}
Under the assumptions above,
\[
\mathrm{BDM2}(X)
\leq
\mathrm{BDM1}(X)
+
C_{\mathrm{rep}}(P^*).
\]
If the baseline representation discrepancies are uniformly bounded and the number of distinct blocks is bounded by the fixed block vocabulary, then
\[
\mathrm{BDM2}(X)
\leq
\mathrm{BDM1}(X)
+
O(1),
\]
where the \(O(1)\) term is independent of \(|X|\), although its hidden constant may be large.
\end{theorem}

\begin{equation}
\label{eq:bdm2-independent-bound}
\mathrm{BDM2}(X)
\leq
\sum_{i=1}^{n}\hat K(p_i^*)
+
\sum_{i=1}^{n}\log m_i.
\end{equation}

\begin{theorem}[Reuse-gain improvement bound]
\label{thm:bdm2-reuse-improvement}
Under the assumptions above, suppose that an admissible BDM\,2.0 explanation uses a set \(R\) of reuse relations. For each \((i,j)\in R\), target \(j\) is described relative to source \(i\), either by recovering \(p_j^*\) from \(p_i^*\) or by recovering \(x_j\) directly from \(x_i\). Each target \(j\) appears at most once in \(R\), and every source index remains outside the target set
\[
T_R
=
\{j:(i,j)\in R\}.
\]
Equivalently, every reuse relation points from a selected representative to a non-selected target.

For any explanatory set \(P=\{p_1,\ldots,p_n\}\in\mathcal E_X\), define the pairwise reuse cost
\[
r_{ij}
=
\min
\left\{
\hat K(p_j\mid p_i),
\hat K(x_j\mid x_i)
\right\},
\]
the corresponding estimator-level reuse gain
\[
\hat{\Delta}_{ij}
=
\hat K(p_j)-r_{ij},
\]
and the total estimator-level reuse gain
\[
\hat{\Delta}
=
\sum_{(i,j)\in R}
\hat{\Delta}_{ij}.
\]
For the shortest-program baseline \(P=P^*\), let
\[
r_{ij}^*
=
\min
\left\{
\hat K(p_j^*\mid p_i^*),
\hat K(x_j\mid x_i)
\right\},
\]
and define
\[
\hat{\Delta}_{ij}^*
=
\hat K(p_j^*)-r_{ij}^*,
\qquad
\hat{\Delta}^*
=
\sum_{(i,j)\in R}
\hat{\Delta}_{ij}^*.
\]
Then
\[
\mathrm{BDM2}(X)
\leq
\mathrm{BDM1}(X)
+
C_{\mathrm{rep}}(P^*)
-
\hat{\Delta}^*.
\]
Consequently, if \(\hat{\Delta}^*>0\), then
\[
\mathrm{BDM2}(X)
<
\mathrm{BDM1}(X)
+
C_{\mathrm{rep}}(P^*).
\]
Moreover, if
\[
\hat{\Delta}^*
>
C_{\mathrm{rep}}(P^*),
\]
then
\[
\mathrm{BDM2}(X)
<
\mathrm{BDM1}(X).
\]
\end{theorem}

\begin{proof}
Let
\[
T_R
=
\{j:(i,j)\in R\}
\]
be the set of reused targets. Consider the choice
\[
P=P^*,
\qquad
S
=
P^*
\setminus
\{p_j^*:j\in T_R\}.
\]
If \(R=\varnothing\), then \(S=P^*\) and the result reduces to Theorem~\ref{thm:bdm2-independent-limit} with \(\hat{\Delta}^*=0\). Otherwise, every reuse relation has a source outside \(T_R\), so \(S\) contains at least one source program. Hence \(S\neq\varnothing\).

For each \((i,j)\in R\), the independent description pays \(\hat K(p_j^*)\), whereas the reused description pays \(r_{ij}^*\). The corresponding estimator-level reuse gain is therefore
\[
\hat{\Delta}_{ij}^*
=
\hat K(p_j^*)-r_{ij}^*.
\]
Summing over all reuse relations gives
\[
\hat{\Delta}^*
=
\sum_{(i,j)\in R}
\hat{\Delta}_{ij}^*.
\]
For the admissible choice above, every source index \(i\) lies outside \(T_R\), so \(p_i^*\in S\). Therefore, for each \((i,j)\in R\), the conditional-description cost of the target \(p_j^*\) satisfies
\[
\min_{p_s^*\in S}
\min
\left\{
\hat K(p_j^*\mid p_s^*),
\hat K(x_j\mid x_s)
\right\}
\leq
\min
\left\{
\hat K(p_j^*\mid p_i^*),
\hat K(x_j\mid x_i)
\right\}
=
r_{ij}^*.
\]
Because \(\mathrm{BDM2}(X)\) minimizes over all admissible explanatory sets and representative subsets, its value is no greater than the cost of this particular reuse construction. Since each target appears at most once in \(R\), summing over the reused targets gives
\[
\mathrm{BDM2}(X)
\leq
\sum_{k\notin T_R}\hat K(p_k^*)
+
\sum_{(i,j)\in R}r_{ij}^*
+
\sum_{k=1}^{n}\log m_k.
\]
By the definition of \(\hat{\Delta}^*\),
\[
\sum_{k\notin T_R}\hat K(p_k^*)
+
\sum_{(i,j)\in R}r_{ij}^*
=
\sum_{k=1}^{n}\hat K(p_k^*)
-
\hat{\Delta}^*.
\]
Hence,
\[
\mathrm{BDM2}(X)
\leq
\sum_{k=1}^{n}\hat K(p_k^*)
+
\sum_{k=1}^{n}\log m_k
-
\hat{\Delta}^*.
\]
This is the independent-description upper bound in Eq.~\eqref{eq:bdm2-independent-bound}, reduced by the total estimator-level reuse gain \(\hat{\Delta}^*\). Applying the same representation-overhead relation used in Theorem~\ref{thm:bdm2-independent-limit},
\[
\hat K(p_k^*)
\leq
\hat K(x_k)
+
c_{\mathrm{rep}}(p_k^*),
\]
we obtain
\[
\begin{aligned}
\mathrm{BDM2}(X)
&\leq
\sum_{k=1}^{n}\hat K(x_k)
+
\sum_{k=1}^{n}\log m_k
+
\sum_{k=1}^{n}c_{\mathrm{rep}}(p_k^*)
-
\hat{\Delta}^* \\
&=
\mathrm{BDM1}(X)
+
C_{\mathrm{rep}}(P^*)
-
\hat{\Delta}^*.
\end{aligned}
\]
Therefore, if \(\hat{\Delta}^*>0\), then
\[
\mathrm{BDM2}(X)
<
\mathrm{BDM1}(X)
+
C_{\mathrm{rep}}(P^*).
\]
Furthermore, if
\[
\hat{\Delta}^*
>
C_{\mathrm{rep}}(P^*),
\]
then
\[
\mathrm{BDM2}(X)
<
\mathrm{BDM1}(X).
\]
\end{proof}

\end{document}
